% ============================================================
%  LEMBAR SOAL  --  FISIKA  --  SSU-ITB 2026  |  Paket 4
%  Sumber: SM-SSU ITB 2025 (30 soal)
% ============================================================

\documentclass[11pt]{article}

\usepackage[utf8]{inputenc}
\usepackage[T1]{fontenc}
\usepackage[a4paper,margin=2.5cm,top=2.8cm]{geometry}
\usepackage{amsmath,amssymb}
\usepackage{graphicx}
\usepackage{enumitem}
\usepackage{multicol}
\usepackage{xcolor}
\usepackage[most]{tcolorbox}
\usepackage{fancyhdr}
\usepackage{booktabs}
\usepackage{icomma}
\usepackage{array}

\definecolor{mapelcolor}{RGB}{20,70,150}
\definecolor{mapelbg}{RGB}{230,238,255}

\pagestyle{fancy}\fancyhf{}
\lhead{\small SSU-ITB 2026}
\rhead{\small Paket 4}
\cfoot{\small\thepage}
\renewcommand{\headrulewidth}{0.4pt}

\newtcolorbox{soal}[1]{
  breakable,
  colback=white, colframe=black!70,
  coltitle=white, colbacktitle=black,
  fonttitle=\bfseries\sffamily,
  title=Nomor #1,
  boxrule=0.8pt, arc=3pt,
  left=3mm, right=3mm, top=2mm, bottom=2mm}

\newenvironment{pil}{%
  \begin{enumerate}[label=(\alph*),leftmargin=*,itemsep=1pt,topsep=3pt]}%
  {\end{enumerate}}
\newenvironment{pilB}{%
  \par\smallskip\begin{multicols}{2}
  \begin{enumerate}[label=(\alph*),leftmargin=*,itemsep=1pt,topsep=3pt]}%
  {\end{enumerate}\end{multicols}}

\newcommand{\gambar}[2][0.52\textwidth]{%
  \begin{center}\includegraphics[width=#1]{#2}\end{center}}

\linespread{1.05}

% ===================================================================
\begin{document}
\thispagestyle{empty}

\begin{center}
  \fbox{\parbox{0.6\textwidth}{\centering\bfseries\large
    LEMBAR SOAL}}
\end{center}
\vspace{0.5cm}

\begin{center}
  {\LARGE\bfseries FISIKA}
\end{center}
\vspace{0.5cm}

\begin{center}
  \begin{tabular}{@{\hspace{2cm}}ll@{\hspace{0.4cm}:@{\hspace{0.4cm}}}l}
    Jenjang      & & SMA / Sederajat               \\[0.3em]
    Hari/Tanggal & & \underline{\hspace{5cm}}      \\[0.3em]
    Waktu        & & \underline{\hspace{5cm}}      \\[0.3em]
    Paket Soal   & & 4                             \\
  \end{tabular}
\end{center}

\vspace{0.6cm}
\noindent\rule{\linewidth}{0.6pt}
\vspace{0.3cm}

\noindent\textbf{\underline{PETUNJUK UMUM}}
\begin{enumerate}[leftmargin=*,itemsep=4pt,topsep=4pt]
  \item Anda \textbf{tidak} diperbolehkan menggunakan sumber-sumber eksternal,
        termasuk internet, \textit{large language model} seperti ChatGPT, Claude,
        LLaMA, Gemini, Meta AI, dan sebagainya. Anda sangat dianjurkan untuk
        mencoba mengerjakan sendiri terlebih dahulu karena evaluasi ini dibuat
        untuk \textbf{mengukur} tingkat kepahaman; nilai $<70$ mengindikasikan
        \textbf{tidak lulus}, dan jika sudah melewati \textit{deadline} maka
        dianggap tidak mengerjakan yang berarti \textbf{tidak lulus}.

  \item Terdapat 2 tahap pengerjaan, yaitu tahap \textbf{Try Out} dan tahap
        \textbf{Review}. Try Out bersifat \textit{open book}; Anda diperbolehkan
        membuka buku apapun selama bersifat \textit{offline} dengan rentang waktu
        2 jam pada tanggal \underline{\hspace{3cm}}. Setelah itu memasuki
        tahap Review dengan \textit{schedule} rentang tanggal
        \underline{\hspace{3cm}}--\underline{\hspace{3cm}}, di mana Anda
        diharapkan mengerjakan ulang dengan disertakan penjelasan/pembelaan
        mengapa revisi jawaban adalah pilihan tersebut, dan tetap mencantumkan
        penjelasan pada soal yang walaupun pada penilaian sudah benar.

  \item Anda dianjurkan untuk mencantumkan caranya walaupun tidak termasuk
        penilaian, mengingat sifatnya adalah sebagai bahan pembelajaran.
        Mohon bertanggung jawab atas pekerjaan Anda sendiri.

  \item Bacalah setiap soal dengan teliti. Soal berbentuk
        \textbf{Pilihan Ganda Biasa} --- 5 opsi (A--E), pilih 1 jawaban
        paling tepat.

  \item Soal berjumlah \textbf{30 butir} soal dengan bobot asli per soal adalah
        $\mathbf{\tfrac{10}{3}}$ untuk mendapatkan nilai sempurna sebesar \textbf{100}.
        Gunakan $g=10$ m/s$^2$ bila tidak disebutkan lain.

  \item Silakan bertanya apabila terdapat soal yang ambigu, rancu, atau kurang
        spesifik selama itu tidak menjurus kepada jawaban soal tersebut.

  \item Isilah detail informasi berikut.\\[0.3em]
        \textbf{Nama}\hspace{0.45cm}:\enspace\underline{\hspace{8cm}}
\end{enumerate}

\vspace{0.6cm}
\noindent\textit{Dengan menandatangani kolom di bawah ini, saya menyatakan telah membaca,
memahami, dan menyetujui seluruh peraturan yang tercantum di atas.}

\vspace{0.6cm}
\noindent\fbox{\parbox{5cm}{\vspace{2.2cm}\centering\small Tanda Tangan Peserta\vspace{0.3cm}}}

\clearpage

% ===================================================================
%  SOAL 1--10
% ===================================================================

\begin{soal}{1 \normalfont\small[PG Biasa]}
\gambar[0.32\textwidth]{images/4/4-n1.png}
Sebuah roda berbentuk silinder pejal memiliki poros di titik $O$. Dua buah
gaya $F_1$ dan $F_2$ bekerja pada roda dan menyebabkan roda berotasi terhadap
porosnya selama 20 detik. Jika setelah itu gaya $F_2$ dihilangkan dan roda
berhenti berotasi setelah 20 detik kemudian, perbandingan $F_2/F_1$ adalah \ldots
\begin{pilB}
  \item $3:2$ \item $2:1$ \item $1:3$ \item $2:3$ \item $1:2$
\end{pilB}
\end{soal}

\begin{soal}{2 \normalfont\small[PG Majemuk Kompleks]}
Sebuah pegas digantung vertikal dengan beban $M$ gram bertambah panjang
$\Delta l$ hingga setimbang. Pegas kemudian berosilasi. Frekuensi osilasi pegas
adalah 11 kali frekuensi bandul bermassa $0{,}5M$ gram panjang tali $L$.
Tentukan kebenaran setiap pernyataan berikut!

\par\smallskip
{\renewcommand{\arraystretch}{1.4}%
\begin{tabular}{@{}p{0.76\linewidth}cc@{}}
\hline\textbf{Pernyataan}&\textbf{Benar}&\textbf{Salah}\\\hline
(1)~Frekuensi pegas $f_s = \dfrac{1}{2\pi}\sqrt{\dfrac{g}{\Delta l}}$ (massa tidak berpengaruh).&$\bigcirc$&$\bigcirc$\\
(2)~Frekuensi bandul $f_b = \dfrac{1}{2\pi}\sqrt{\dfrac{g}{L}}$ tidak bergantung massa beban.&$\bigcirc$&$\bigcirc$\\
(3)~Dari syarat $f_s = 11f_b$, diperoleh $L = 11\,\Delta l$.&$\bigcirc$&$\bigcirc$\\
(4)~Jika $\Delta l = 1{,}9$ cm, maka $L = 229{,}9$ cm.&$\bigcirc$&$\bigcirc$\\\hline
\end{tabular}}
\end{soal}

\begin{soal}{3 \normalfont\small[PG Biasa]}
Persamaan gelombang sinusoidal yang merambat ke arah negatif $x$ dengan
amplitudo $10$ cm, panjang gelombang $314$ cm, frekuensi $2$ Hz, dan
$y(x{=}10\text{ cm},\,t{=}0)=0$ adalah \ldots
\begin{pil}
  \item $y=0{,}1\sin(2x-4\pi t-0{,}2)$
  \item $y=0{,}1\sin(2x+4\pi t)$
  \item $y=0{,}1\sin(2x+4\pi t-0{,}2)$
  \item $y=0{,}1\sin(2x+4\pi t+0{,}2)$
  \item $y=0{,}1\sin(2x-4\pi t)$
\end{pil}
\end{soal}

\begin{soal}{4 \normalfont\small[PG Biasa]}
\gambar[0.3\textwidth]{images/4/4-2.png}
Balok $A$ bermassa $M$ digantung tali panjang $L$, disimpangkan hingga
ketinggian $H$ lalu dilepas. Di titik terendah, $A$ menumbuk benda $B$
bermassa $m$ yang diam di permukaan licin dan keduanya menempel. Ketinggian
maksimum sistem setelah tumbukan adalah \ldots $H$.
\begin{pilB}
  \item $\left(\dfrac{M+m}{m}\right)^{\!2}$
  \item $\left(\dfrac{M}{M+m}\right)^{\!2}$
  \item $\left(\dfrac{m}{M+m}\right)^{\!2}$
  \item $\left(\dfrac{M}{M+m}\right)$
  \item $\left(\dfrac{M+m}{M}\right)^{\!2}$
\end{pilB}
\end{soal}

\begin{soal}{5 \normalfont\small[PG Biasa]}
Apa pengaruh angin terhadap efek Doppler pada suara yang merambat di udara?
\begin{pil}
  \item Efek Doppler terjadi selama ada angin meskipun sumber dan pengamat diam.
  \item Angin dapat mengubah besar efek Doppler, tetapi tidak dapat menyebabkannya.
  \item Angin dapat menyebabkan efek Doppler, seperti halnya sumber yang bergerak.
  \item Angin tidak memengaruhi efek Doppler sama sekali.
  \item Efek Doppler hanya terjadi jika angin bertiup ke arah pengamat.
\end{pil}
\end{soal}

\begin{soal}{6 \normalfont\small[PG Majemuk Kompleks]}
Sebuah balok berada di bidang datar kasar. Balok ditarik dengan gaya arah
$0^\circ<\alpha<80^\circ$ terhadap perpindahan sehingga bergerak.
Tentukan kebenaran setiap pernyataan berikut!

\par\smallskip
{\renewcommand{\arraystretch}{1.4}%
\begin{tabular}{@{}p{0.76\linewidth}cc@{}}
\hline\textbf{Pernyataan}&\textbf{Benar}&\textbf{Salah}\\\hline
(1)~Kerja oleh gaya gesek pada balok bertanda positif.&$\bigcirc$&$\bigcirc$\\
(2)~Peningkatan energi kinetik balok sebanding dengan perpindahan.&$\bigcirc$&$\bigcirc$\\
(3)~Jika $\alpha$ berubah, hanya kerja dari gaya tarik $F$ yang terpengaruh.&$\bigcirc$&$\bigcirc$\\
(4)~Kerja gaya gravitasi dan gaya normal masing-masing bernilai nol.&$\bigcirc$&$\bigcirc$\\\hline
\end{tabular}}
\end{soal}

\begin{soal}{7 \normalfont\small[PG Majemuk Kompleks]}
Gelombang berdiri dengan ujung terikat:
$y=0{,}2\sin(5\pi x)\cos(2\pi t)$, $x$ dan $y$ dalam cm, $t$ dalam sekon.
Tentukan kebenaran setiap pernyataan berikut!

\par\smallskip
{\renewcommand{\arraystretch}{1.4}%
\begin{tabular}{@{}p{0.76\linewidth}cc@{}}
\hline\textbf{Pernyataan}&\textbf{Benar}&\textbf{Salah}\\\hline
(1)~Periode gelombang adalah 1 sekon.&$\bigcirc$&$\bigcirc$\\
(2)~Amplitudo gelombang bernilai konstan, yaitu 0{,}2 cm.&$\bigcirc$&$\bigcirc$\\
(3)~Pada jarak 0{,}4 cm dari ujung terikat terbentuk simpul.&$\bigcirc$&$\bigcirc$\\
(4)~Cepat rambat gelombang adalah 0{,}4 m/s.&$\bigcirc$&$\bigcirc$\\\hline
\end{tabular}}
\end{soal}

\begin{soal}{8 \normalfont\small[PG Biasa]}
\gambar[0.42\textwidth]{images/4/4-n8.png}
Balok $A$ berat 100 N berada di atas balok $B$ dan diikat tali mendatar.
Balok $B$ berat 500 N. Koefisien gesek statik dan kinetik antara $A$--$B$:
0{,}2 dan 0{,}1; antara $B$--lantai: 0{,}6 dan 0{,}5. Gaya mendatar $F$
minimum untuk menggeser balok $B$ adalah \ldots\ N.
\begin{pilB}
  \item $380$ \item $400$ \item $420$ \item $440$ \item $460$
\end{pilB}
\end{soal}

\begin{soal}{9 \normalfont\small[PG Biasa]}
\gambar[0.4\textwidth]{images/4/4-n9.png}
Dua buah vektor $A$ dan $B$ terletak pada bidang $xy$ seperti pada gambar.
Jika komponen resultan kedua vektor pada sumbu $x$ sama dengan nol,
besar vektor $A$ adalah \ldots\ N.
\begin{pilB}
  \item $6$ \item $8$ \item $10$ \item $12$ \item $14$
\end{pilB}
\end{soal}

\begin{soal}{10 \normalfont\small[PG Majemuk Kompleks]}
Posisi suatu partikel: $x(t)=6t^2-4t+2$ (m, s).
Tentukan kebenaran setiap pernyataan berikut!

\par\smallskip
{\renewcommand{\arraystretch}{1.4}%
\begin{tabular}{@{}p{0.76\linewidth}cc@{}}
\hline\textbf{Pernyataan}&\textbf{Benar}&\textbf{Salah}\\\hline
(1)~Kecepatan partikel $v(t)=12t+4$ m/s.&$\bigcirc$&$\bigcirc$\\
(2)~Percepatan partikel $a=12$ m/s$^2$ (konstan).&$\bigcirc$&$\bigcirc$\\
(3)~Kecepatan pada $t=3$ s adalah $v(3)=32$ m/s.&$\bigcirc$&$\bigcirc$\\
(4)~Kecepatan awal $v(0)=-4$ m/s.&$\bigcirc$&$\bigcirc$\\\hline
\end{tabular}}
\end{soal}

% ===================================================================
%  SOAL 11--20
% ===================================================================

\begin{soal}{11 \normalfont\small[PG Biasa]}
Partikel $Z$ memiliki energi total 17{,}0 GeV dan massa diam
$8{,}0$ GeV/$c^2$. Energi kinetiknya adalah \ldots\ GeV.
\begin{pilB}
  \item $17{,}0$ \item $25{,}0$ \item $136{,}0$ \item $8{,}0$ \item $9{,}0$
\end{pilB}
\end{soal}

\begin{soal}{12 \normalfont\small[PG Biasa]}
Dua bola identik dilempar horizontal ke dinding vertikal licin dengan kelajuan
sama. Bola $P$ memantul dengan kelajuan sama; bola $Q$ berhenti. Perbandingan
impuls $I_P/I_Q$ dan perbandingan gaya rata-rata $F_P/F_Q$ adalah \ldots
\begin{pil}
  \item $I_P/I_Q=1$ dan $F_P/F_Q=1$
  \item $I_P/I_Q=1$ dan $F_P/F_Q=2$
  \item $I_P/I_Q=2$ dan $F_P/F_Q=1$
  \item $I_P/I_Q=2$ dan $F_P/F_Q=2$
  \item $I_P/I_Q=2$ dan $F_P/F_Q=\Delta t_Q/\Delta t_P$
\end{pil}
\end{soal}

\begin{soal}{13 \normalfont\small[PG Biasa]}
\gambar[0.26\textwidth]{images/4/4-3.png}
Bola karet massa 0{,}20 kg dijatuhkan dari ketinggian 1{,}8 m di atas pegas
vertikal. Saat kompresi maksimum, gaya pegas terbesar yang aman adalah 36 N.
Abaikan hambatan udara, $g=10$ m/s$^2$. Konstanta pegas minimum agar syarat
terpenuhi adalah \ldots\ N/m.
\begin{pilB}
  \item $80$ \item $120$ \item $160$ \item $200$ \item $240$
\end{pilB}
\end{soal}

\begin{soal}{14 \normalfont\small[PG Majemuk Kompleks]}
\gambar[0.32\textwidth]{images/4/4-4.png}
Bola $P$ massa 0{,}30 kg diangkat hingga ketinggian 20 cm, dilepas, lalu
menumbuk bola $Q$ massa 0{,}20 kg yang diam di titik terendah; keduanya
menempel. Tentukan kebenaran setiap pernyataan berikut!

\par\smallskip
{\renewcommand{\arraystretch}{1.4}%
\begin{tabular}{@{}p{0.76\linewidth}cc@{}}
\hline\textbf{Pernyataan}&\textbf{Benar}&\textbf{Salah}\\\hline
(1)~Kecepatan $P$ di titik terendah sebelum tumbukan: $v=2$ m/s.&$\bigcirc$&$\bigcirc$\\
(2)~Setelah tumbukan menempel, $v'=0{,}6v=1{,}2$ m/s.&$\bigcirc$&$\bigcirc$\\
(3)~Tinggi maksimum setelah tumbukan: $h'=0{,}6H=12$ cm.&$\bigcirc$&$\bigcirc$\\
(4)~Tinggi maksimum setelah tumbukan: $h'=0{,}36H=7{,}2$ cm.&$\bigcirc$&$\bigcirc$\\\hline
\end{tabular}}
\end{soal}

\begin{soal}{15 \normalfont\small[PG Biasa]}
Suatu planet $X$ memiliki jari-jari 1{,}5 kali jari-jari Bumi dan massa jenis
rata-rata 2 kali massa jenis Bumi. Jika berat suatu benda di Bumi adalah $W$,
berat benda tersebut di planet $X$ adalah \ldots
\begin{pilB}
  \item $\tfrac{2}{3}W$ \item $W$ \item $\tfrac{3}{2}W$ \item $2W$ \item $3W$
\end{pilB}
\end{soal}

\begin{soal}{16 \normalfont\small[Isian Singkat]}
\gambar[0.5\textwidth]{images/4/4-5.png}
Batang homogen panjang 4 m berat 100 N ditumpu di kedua ujungnya.
Anak berat 300 N berdiri 1 m dari ujung kiri; beban 200 N digantung
0{,}5 m dari ujung kanan. Besar gaya tumpu pada ujung kanan adalah \ldots\ N.

\vspace{0.4em}
\noindent Jawaban:\enspace\underline{\hspace{4cm}}
\end{soal}

\begin{soal}{17 \normalfont\small[PG Majemuk Kompleks]}
\gambar[0.42\textwidth]{images/4/4-6.png}
Dua pegas identik masing-masing berkonstanta $k$ disusun seri. Saat ditarik
20 cm dari setimbang, usaha yang diperlukan 4 J.
Tentukan kebenaran setiap pernyataan berikut!

\par\smallskip
{\renewcommand{\arraystretch}{1.4}%
\begin{tabular}{@{}p{0.76\linewidth}cc@{}}
\hline\textbf{Pernyataan}&\textbf{Benar}&\textbf{Salah}\\\hline
(1)~Konstanta ekivalen dua pegas seri identik: $k_{\text{eq}}=\dfrac{k}{2}$.&$\bigcirc$&$\bigcirc$\\
(2)~Dari data $W=4$ J dan $x=0{,}20$ m, diperoleh $k=200$ N/m.&$\bigcirc$&$\bigcirc$\\
(3)~Nilai $k$ untuk setiap pegas adalah 400 N/m.&$\bigcirc$&$\bigcirc$\\
(4)~Jika disusun paralel, konstanta ekivalen kedua pegas adalah 800 N/m.&$\bigcirc$&$\bigcirc$\\\hline
\end{tabular}}
\end{soal}

\begin{soal}{18 \normalfont\small[Isian Singkat]}
\gambar[0.42\textwidth]{images/4/4-n18.png}
Gaya interaksi selama tumbukan membentuk grafik segitiga $F$-$t$ dengan
puncak 160 N dan lama tumbukan 0{,}010 s. Bola massa 0{,}05 kg bergerak
mendekati dinding dengan kelajuan 12 m/s, lalu berbalik arah. Besar kelajuan
sesudah tumbukan adalah \ldots\ m/s.

\vspace{0.4em}
\noindent Jawaban:\enspace\underline{\hspace{4cm}}
\end{soal}

\begin{soal}{19 \normalfont\small[PG Majemuk Kompleks]}
Energi potensial pegas $U=\tfrac{1}{2}kx^2$. Pada simpangan $x=10$ cm,
$U=0{,}80$ J. Benda massa 0{,}50 kg berosilasi pada pegas itu.
Tentukan kebenaran setiap pernyataan berikut!

\par\smallskip
{\renewcommand{\arraystretch}{1.4}%
\begin{tabular}{@{}p{0.76\linewidth}cc@{}}
\hline\textbf{Pernyataan}&\textbf{Benar}&\textbf{Salah}\\\hline
(1)~Dari data: $k=\dfrac{2U}{x^2}=160$ N/m.&$\bigcirc$&$\bigcirc$\\
(2)~Periode osilasi massa-pegas: $T=2\pi\sqrt{\dfrac{m}{k}}$.&$\bigcirc$&$\bigcirc$\\
(3)~Dengan $m=0{,}50$ kg dan $k=160$ N/m, periode $T\approx0{,}50$ s.&$\bigcirc$&$\bigcirc$\\
(4)~Dengan $m=0{,}50$ kg dan $k=160$ N/m, periode $T\approx0{,}35$ s.&$\bigcirc$&$\bigcirc$\\\hline
\end{tabular}}
\end{soal}

\begin{soal}{20 \normalfont\small[PG Majemuk Kompleks]}
Dua getaran harmonik segaris: $y_1=4\sin(10t)$ cm dan $y_2=3\cos(10t)$ cm.
Tentukan kebenaran setiap pernyataan berikut!

\par\smallskip
{\renewcommand{\arraystretch}{1.4}%
\begin{tabular}{@{}p{0.76\linewidth}cc@{}}
\hline\textbf{Pernyataan}&\textbf{Benar}&\textbf{Salah}\\\hline
(1)~Amplitudo resultan kedua getaran adalah $4+3=7$ cm.&$\bigcirc$&$\bigcirc$\\
(2)~$y_2=3\cos(10t)=3\sin(10t+90^\circ)$; beda fase $y_1$ dan $y_2$ adalah $90^\circ$.&$\bigcirc$&$\bigcirc$\\
(3)~Amplitudo resultan $R=\sqrt{4^2+3^2}=5$ cm.&$\bigcirc$&$\bigcirc$\\
(4)~Kedua getaran memiliki frekuensi sudut yang sama, yaitu 10 rad/s.&$\bigcirc$&$\bigcirc$\\\hline
\end{tabular}}
\end{soal}

% ===================================================================
%  SOAL 21--30
% ===================================================================

\begin{soal}{21 \normalfont\small[PG Biasa]}
\gambar[0.4\textwidth]{images/4/4-7.png}
Partikel melakukan GHS amplitudo 6 cm, frekuensi 0{,}5 Hz. Pada $t=0$
partikel di titik setimbang bergerak ke kiri. Posisi partikel 0{,}5 s
kemudian adalah $x=\ldots$ cm.
\begin{pilB}
  \item $-6$ \item $-3$ \item $0$ \item $3$ \item $6$
\end{pilB}
\end{soal}

\begin{soal}{22 \normalfont\small[Isian Singkat]}
Pipa organa terbuka panjang 85 cm. Cepat rambat bunyi di udara 340 m/s.
Frekuensi nada atas kedua yang dapat dihasilkan pipa tersebut adalah \ldots\ Hz.

\vspace{0.4em}
\noindent Jawaban:\enspace\underline{\hspace{4cm}}
\end{soal}

\begin{soal}{23 \normalfont\small[Isian Singkat]}
Gelombang pada tali: $y=0{,}04\sin(\pi x-20\pi t)$ ($x$, $y$ dalam m; $t$
dalam s). Massa tali 0{,}80 g, panjang 2{,}0 m. Tegangan tali adalah \ldots\ N.

\vspace{0.4em}
\noindent Jawaban:\enspace\underline{\hspace{4cm}}
\end{soal}

\begin{soal}{24 \normalfont\small[PG Majemuk Kompleks]}
Gelombang transversal merambat ke kiri: amplitudo 5 cm, panjang gelombang
2 m, frekuensi 4 Hz. Pada $x=0$ dan $t=0$: simpangan nol, medium bergerak
ke atas. Tentukan kebenaran setiap pernyataan berikut!

\par\smallskip
{\renewcommand{\arraystretch}{1.4}%
\begin{tabular}{@{}p{0.76\linewidth}cc@{}}
\hline\textbf{Pernyataan}&\textbf{Benar}&\textbf{Salah}\\\hline
(1)~$k=\pi$ rad/m dan $\omega=8\pi$ rad/s.&$\bigcirc$&$\bigcirc$\\
(2)~Gelombang ke kiri: tanda positif pada $\omega t$ dalam argumen sinus.&$\bigcirc$&$\bigcirc$\\
(3)~Pada $x=0$, $t=0$: simpangan $y=0{,}05\sin(0)=0$ sesuai syarat.&$\bigcirc$&$\bigcirc$\\
(4)~Pada $x=0$, $t=0$, medium bergerak ke bawah (kecepatan partikel negatif).&$\bigcirc$&$\bigcirc$\\\hline
\end{tabular}}
\end{soal}

\begin{soal}{25 \normalfont\small[Isian Singkat]}
Gelombang stasioner pada tali: $y=10\sin(4\pi x)\cos(50\pi t)$,
$x$ dalam m, $y$ dalam cm. Jarak antara dua simpul yang berurutan
adalah \ldots\ m.

\vspace{0.4em}
\noindent Jawaban:\enspace\underline{\hspace{4cm}}
\end{soal}

\begin{soal}{26 \normalfont\small[PG Majemuk Kompleks]}
\gambar[0.5\textwidth]{images/4/4-8.png}
Dalam foto sesaat gelombang pada tali, jarak puncak pertama ke puncak ketiga
adalah 1{,}2 m. Dalam 4 s, sumber menghasilkan 8 gelombang penuh.
Tentukan kebenaran setiap pernyataan berikut!

\par\smallskip
{\renewcommand{\arraystretch}{1.4}%
\begin{tabular}{@{}p{0.76\linewidth}cc@{}}
\hline\textbf{Pernyataan}&\textbf{Benar}&\textbf{Salah}\\\hline
(1)~Jarak puncak pertama ke puncak ketiga mencakup dua panjang gelombang.&$\bigcirc$&$\bigcirc$\\
(2)~Dari data: $2\lambda=1{,}2$ m, sehingga $\lambda=0{,}6$ m.&$\bigcirc$&$\bigcirc$\\
(3)~8 gelombang dalam 4 s memberikan frekuensi $f=2$ Hz.&$\bigcirc$&$\bigcirc$\\
(4)~Cepat rambat gelombang $v=f\lambda=2\times0{,}6=0{,}6$ m/s.&$\bigcirc$&$\bigcirc$\\\hline
\end{tabular}}
\end{soal}

\begin{soal}{27 \normalfont\small[PG Majemuk Kompleks]}
$y_1=3\cos(kx-\omega t)$ dan $y_2=3\cos(kx+\omega t)$.
Tentukan kebenaran setiap pernyataan berikut!

\par\smallskip
{\renewcommand{\arraystretch}{1.4}%
\begin{tabular}{@{}p{0.76\linewidth}cc@{}}
\hline\textbf{Pernyataan}&\textbf{Benar}&\textbf{Salah}\\\hline
(1)~Resultan $y_1+y_2$ merupakan gelombang stasioner.&$\bigcirc$&$\bigcirc$\\
(2)~Amplitudo maksimum resultan adalah 6.&$\bigcirc$&$\bigcirc$\\
(3)~Simpul terjadi saat $kx=\dfrac{(2n+1)\pi}{2}$, $n\in\mathbb{Z}$.&$\bigcirc$&$\bigcirc$\\
(4)~Kedua gelombang memiliki frekuensi yang sama.&$\bigcirc$&$\bigcirc$\\\hline
\end{tabular}}
\end{soal}

\begin{soal}{28 \normalfont\small[PG Biasa]}
Dua gelombang frekuensi dan panjang gelombang sama merambat berlawanan arah:
$y_1=4\cos(kx-\omega t)$ dan $y_2=6\cos(kx+\omega t)$. Amplitudo maksimum
dan amplitudo minimum resultan berturut-turut adalah \ldots
\begin{pilB}
  \item $2$ dan $10$ \item $4$ dan $6$ \item $6$ dan $4$
  \item $10$ dan $2$ \item $12$ dan $0$
\end{pilB}
\end{soal}

\begin{soal}{29 \normalfont\small[Isian Singkat]}
Jendela terbuka luas 0{,}40 m$^2$ dilewati bunyi taraf intensitas 80 dB.
($I_0=10^{-12}$ W/m$^2$.) Daya akustik yang masuk adalah $4\times\ldots$ W.

\vspace{0.4em}
\noindent Jawaban:\enspace\underline{\hspace{4cm}}
\end{soal}

\begin{soal}{30 \normalfont\small[PG Majemuk Kompleks]}
Speaker titik memancarkan daya tetap ke segala arah. Taraf intensitas pada
jarak 2 m adalah 70 dB. Tentukan kebenaran setiap pernyataan berikut!

\par\smallskip
{\renewcommand{\arraystretch}{1.4}%
\begin{tabular}{@{}p{0.76\linewidth}cc@{}}
\hline\textbf{Pernyataan}&\textbf{Benar}&\textbf{Salah}\\\hline
(1)~Untuk sumber titik, $I\propto\dfrac{1}{r^2}$.&$\bigcirc$&$\bigcirc$\\
(2)~Saat jarak bertambah 10 kali, intensitas berkurang 10 kali ($\Delta\beta=-10$ dB).&$\bigcirc$&$\bigcirc$\\
(3)~Perubahan taraf: $\Delta\beta=10\log(1/100)=-20$ dB.&$\bigcirc$&$\bigcirc$\\
(4)~Taraf intensitas pada jarak 20 m: $\beta_2=50$ dB.&$\bigcirc$&$\bigcirc$\\\hline
\end{tabular}}
\end{soal}

\end{document}
